\section{Banach algebras}


\begin{exercise}{1}
Why is the vector space $X$ of all complex $n \times n$ matrices ($n > 1$, fixed) with identity $I$ (the $n$-rowed unit matrix) complete?
\end{exercise}
\begin{proof}
We can identify $X$ with $\C^{n^2}$, which we know is complete from 1.5-1, thus $X$ is itself complete.
Since in finite dimension every norm is equivalent, then the choice of norm does not affect this result.
\end{proof} 

\begin{exercise}{2}
Show that $\norm{x y} \leq \norm{x} \norm{y}$ holds in the space $C[a,b]$ with identity $e = 1$ the product being defined as usual $(xy)(t) = x(t)y(t)$.
\end{exercise}
\begin{proof}
Note,
\begin{align*}
    \norm{x y}
    = \max_t \absoluteValue{x(t) y(t)}
    \leq \max_t \absoluteValue{x(t)} \max_t \absoluteValue{y(t)}
    = \norm{x}\norm{y},
\end{align*}
as required.
\end{proof} 

\begin{exercise}{3}
How can we make the vector space of all ordered $n$-tuples of complex numbers into a Banach algebra?
\end{exercise}
\begin{proof}
We just need to define a suitable multiplication, which satisfies all the properties of the Algebra and a norm under which it will be a complete space.
For the product, elementwise multiplication suffices, and for the norm, the max norm has the property that $\norm{x y} \leq \norm{x} \norm{y}$ for all $x, y \in \C^n$.
\end{proof} 

\begin{exercise}{4}
What are the invertible elements in the spaces of exercises 1, 2, and $X$ given by $\R$, $\C$, with identity $e=1$?
\end{exercise}
\begin{proof}
In the exercise 1, the set of invertible matrices, or matrices with determinant different from 0.
In exercise 2, the functions which do not evaluate to 0, since the space considered is $C[a,b]$, then these functions are the strictly positive or strictly negative functions.
In $\R$ and $\C$ any element with the exception of 0.
\end{proof} 

\begin{exercise}{5}
Show that for the elements of $X$ in exercise 1, the definition of a spectrum in 7.6-6 agrees with the definition of the spectrum in finite dimensional spaces (7.1-1).
\end{exercise}
\begin{proof}
To see this agreement, we just need to translate one definition to the other.
Since the set of $n$ by $n$ complex matrices forms a Banach algebra, the spectrum of any matrix, $x$, is defined as the set of scalars such that $(x - \lambda e)^{-1}$ is not defined.
In the case of complex matrices, $e = I$, and so this is precisely the definition of eigenvalues and point spectrum (which in finite dimensions corresponds to the whole spectrum) of matrices.
\end{proof} 

\begin{exercise}{6}
Find $\sigma(x)$ of $x \in C[0, 2\pi]$, where $x(t) = \sin t$.
Find $\sigma(x)$ for any $x \in C[a,b]$.
\end{exercise}
\begin{proof}
For any $x \in C[a,b]$ (including $x(t) = \sin t$ in $C[0,2\pi]$), $\sigma(x)$ is simply the of $x$, since for $\lambda \in \cR(x)$, $(x - \lambda e)(t) = 0$, so it is not invertible, and thus $\lambda \in \sigma(x)$.
\end{proof} 

\begin{exercise}{8}
Let $A$ be a complex Banach algebra with identity $e$.
If for an $x \in A$ there are $y, z \in A$ such that $yx = e$ and $xz = e$, show that $x$ is invertible and $y = z = x^{-1}$.
\end{exercise}
\begin{proof}
We have $x z x = x$, and $y x z x = y x$, simplifying $yx = e$ on both sides of the equality, we get $z x = e$, since $z x = e = x z$, then $x$ is invertible with $z = x^{-1}$.
We can reason in a similar way to conclude that $x^{-1} = y$.
Since inverses are unique, then $y = z$, completing the proof.
\end{proof} 

\begin{exercise}{9}
If $x \in A$ is invertible and commutes with $y \in A$, show that $x^{-1}$ and $y$ also commute.
\end{exercise}
\begin{proof}
We have $x^{-1} y = x^{-1} y x x^{-1} = x^{-1} x y x^{-1} = y x^{-1}$, as required.
\end{proof} 
